% Options for packages loaded elsewhere
% Options for packages loaded elsewhere
\PassOptionsToPackage{unicode}{hyperref}
\PassOptionsToPackage{hyphens}{url}
%
\documentclass[
  letterpaper,
]{book}
\usepackage{xcolor}
\usepackage{amsmath,amssymb}
\setcounter{secnumdepth}{5}
\usepackage{iftex}
\ifPDFTeX
  \usepackage[T1]{fontenc}
  \usepackage[utf8]{inputenc}
  \usepackage{textcomp} % provide euro and other symbols
\else % if luatex or xetex
  \usepackage{unicode-math} % this also loads fontspec
  \defaultfontfeatures{Scale=MatchLowercase}
  \defaultfontfeatures[\rmfamily]{Ligatures=TeX,Scale=1}
\fi
\usepackage{lmodern}
\ifPDFTeX\else
  % xetex/luatex font selection
\fi
% Use upquote if available, for straight quotes in verbatim environments
\IfFileExists{upquote.sty}{\usepackage{upquote}}{}
\IfFileExists{microtype.sty}{% use microtype if available
  \usepackage[]{microtype}
  \UseMicrotypeSet[protrusion]{basicmath} % disable protrusion for tt fonts
}{}
\makeatletter
\@ifundefined{KOMAClassName}{% if non-KOMA class
  \IfFileExists{parskip.sty}{%
    \usepackage{parskip}
  }{% else
    \setlength{\parindent}{0pt}
    \setlength{\parskip}{6pt plus 2pt minus 1pt}}
}{% if KOMA class
  \KOMAoptions{parskip=half}}
\makeatother
% Make \paragraph and \subparagraph free-standing
\makeatletter
\ifx\paragraph\undefined\else
  \let\oldparagraph\paragraph
  \renewcommand{\paragraph}{
    \@ifstar
      \xxxParagraphStar
      \xxxParagraphNoStar
  }
  \newcommand{\xxxParagraphStar}[1]{\oldparagraph*{#1}\mbox{}}
  \newcommand{\xxxParagraphNoStar}[1]{\oldparagraph{#1}\mbox{}}
\fi
\ifx\subparagraph\undefined\else
  \let\oldsubparagraph\subparagraph
  \renewcommand{\subparagraph}{
    \@ifstar
      \xxxSubParagraphStar
      \xxxSubParagraphNoStar
  }
  \newcommand{\xxxSubParagraphStar}[1]{\oldsubparagraph*{#1}\mbox{}}
  \newcommand{\xxxSubParagraphNoStar}[1]{\oldsubparagraph{#1}\mbox{}}
\fi
\makeatother

\usepackage{color}
\usepackage{fancyvrb}
\newcommand{\VerbBar}{|}
\newcommand{\VERB}{\Verb[commandchars=\\\{\}]}
\DefineVerbatimEnvironment{Highlighting}{Verbatim}{commandchars=\\\{\}}
% Add ',fontsize=\small' for more characters per line
\usepackage{framed}
\definecolor{shadecolor}{RGB}{241,243,245}
\newenvironment{Shaded}{\begin{snugshade}}{\end{snugshade}}
\newcommand{\AlertTok}[1]{\textcolor[rgb]{0.68,0.00,0.00}{#1}}
\newcommand{\AnnotationTok}[1]{\textcolor[rgb]{0.37,0.37,0.37}{#1}}
\newcommand{\AttributeTok}[1]{\textcolor[rgb]{0.40,0.45,0.13}{#1}}
\newcommand{\BaseNTok}[1]{\textcolor[rgb]{0.68,0.00,0.00}{#1}}
\newcommand{\BuiltInTok}[1]{\textcolor[rgb]{0.00,0.23,0.31}{#1}}
\newcommand{\CharTok}[1]{\textcolor[rgb]{0.13,0.47,0.30}{#1}}
\newcommand{\CommentTok}[1]{\textcolor[rgb]{0.37,0.37,0.37}{#1}}
\newcommand{\CommentVarTok}[1]{\textcolor[rgb]{0.37,0.37,0.37}{\textit{#1}}}
\newcommand{\ConstantTok}[1]{\textcolor[rgb]{0.56,0.35,0.01}{#1}}
\newcommand{\ControlFlowTok}[1]{\textcolor[rgb]{0.00,0.23,0.31}{\textbf{#1}}}
\newcommand{\DataTypeTok}[1]{\textcolor[rgb]{0.68,0.00,0.00}{#1}}
\newcommand{\DecValTok}[1]{\textcolor[rgb]{0.68,0.00,0.00}{#1}}
\newcommand{\DocumentationTok}[1]{\textcolor[rgb]{0.37,0.37,0.37}{\textit{#1}}}
\newcommand{\ErrorTok}[1]{\textcolor[rgb]{0.68,0.00,0.00}{#1}}
\newcommand{\ExtensionTok}[1]{\textcolor[rgb]{0.00,0.23,0.31}{#1}}
\newcommand{\FloatTok}[1]{\textcolor[rgb]{0.68,0.00,0.00}{#1}}
\newcommand{\FunctionTok}[1]{\textcolor[rgb]{0.28,0.35,0.67}{#1}}
\newcommand{\ImportTok}[1]{\textcolor[rgb]{0.00,0.46,0.62}{#1}}
\newcommand{\InformationTok}[1]{\textcolor[rgb]{0.37,0.37,0.37}{#1}}
\newcommand{\KeywordTok}[1]{\textcolor[rgb]{0.00,0.23,0.31}{\textbf{#1}}}
\newcommand{\NormalTok}[1]{\textcolor[rgb]{0.00,0.23,0.31}{#1}}
\newcommand{\OperatorTok}[1]{\textcolor[rgb]{0.37,0.37,0.37}{#1}}
\newcommand{\OtherTok}[1]{\textcolor[rgb]{0.00,0.23,0.31}{#1}}
\newcommand{\PreprocessorTok}[1]{\textcolor[rgb]{0.68,0.00,0.00}{#1}}
\newcommand{\RegionMarkerTok}[1]{\textcolor[rgb]{0.00,0.23,0.31}{#1}}
\newcommand{\SpecialCharTok}[1]{\textcolor[rgb]{0.37,0.37,0.37}{#1}}
\newcommand{\SpecialStringTok}[1]{\textcolor[rgb]{0.13,0.47,0.30}{#1}}
\newcommand{\StringTok}[1]{\textcolor[rgb]{0.13,0.47,0.30}{#1}}
\newcommand{\VariableTok}[1]{\textcolor[rgb]{0.07,0.07,0.07}{#1}}
\newcommand{\VerbatimStringTok}[1]{\textcolor[rgb]{0.13,0.47,0.30}{#1}}
\newcommand{\WarningTok}[1]{\textcolor[rgb]{0.37,0.37,0.37}{\textit{#1}}}

\usepackage{longtable,booktabs,array}
\newcounter{none} % for unnumbered tables
\usepackage{calc} % for calculating minipage widths
% Correct order of tables after \paragraph or \subparagraph
\usepackage{etoolbox}
\makeatletter
\patchcmd\longtable{\par}{\if@noskipsec\mbox{}\fi\par}{}{}
\makeatother
% Allow footnotes in longtable head/foot
\IfFileExists{footnotehyper.sty}{\usepackage{footnotehyper}}{\usepackage{footnote}}
\makesavenoteenv{longtable}
\usepackage{graphicx}
\makeatletter
\newsavebox\pandoc@box
\newcommand*\pandocbounded[1]{% scales image to fit in text height/width
  \sbox\pandoc@box{#1}%
  \Gscale@div\@tempa{\textheight}{\dimexpr\ht\pandoc@box+\dp\pandoc@box\relax}%
  \Gscale@div\@tempb{\linewidth}{\wd\pandoc@box}%
  \ifdim\@tempb\p@<\@tempa\p@\let\@tempa\@tempb\fi% select the smaller of both
  \ifdim\@tempa\p@<\p@\scalebox{\@tempa}{\usebox\pandoc@box}%
  \else\usebox{\pandoc@box}%
  \fi%
}
% Set default figure placement to htbp
\def\fps@figure{htbp}
\makeatother





\setlength{\emergencystretch}{3em} % prevent overfull lines

\providecommand{\tightlist}{%
  \setlength{\itemsep}{0pt}\setlength{\parskip}{0pt}}



 


\makeatletter
\@ifpackageloaded{bookmark}{}{\usepackage{bookmark}}
\makeatother
\makeatletter
\@ifpackageloaded{caption}{}{\usepackage{caption}}
\AtBeginDocument{%
\ifdefined\contentsname
  \renewcommand*\contentsname{Table of contents}
\else
  \newcommand\contentsname{Table of contents}
\fi
\ifdefined\listfigurename
  \renewcommand*\listfigurename{List of Figures}
\else
  \newcommand\listfigurename{List of Figures}
\fi
\ifdefined\listtablename
  \renewcommand*\listtablename{List of Tables}
\else
  \newcommand\listtablename{List of Tables}
\fi
\ifdefined\figurename
  \renewcommand*\figurename{Figure}
\else
  \newcommand\figurename{Figure}
\fi
\ifdefined\tablename
  \renewcommand*\tablename{Table}
\else
  \newcommand\tablename{Table}
\fi
}
\@ifpackageloaded{float}{}{\usepackage{float}}
\floatstyle{ruled}
\@ifundefined{c@chapter}{\newfloat{codelisting}{h}{lop}}{\newfloat{codelisting}{h}{lop}[chapter]}
\floatname{codelisting}{Listing}
\newcommand*\listoflistings{\listof{codelisting}{List of Listings}}
\makeatother
\makeatletter
\makeatother
\makeatletter
\@ifpackageloaded{caption}{}{\usepackage{caption}}
\@ifpackageloaded{subcaption}{}{\usepackage{subcaption}}
\makeatother
\usepackage{bookmark}
\IfFileExists{xurl.sty}{\usepackage{xurl}}{} % add URL line breaks if available
\urlstyle{same}
\hypersetup{
  pdftitle={Environmental Informatics Using Research Infrastructures and their Data},
  pdfauthor={Dr.~Katharyn Duffy, Dr.~Ben Ruddell},
  hidelinks,
  pdfcreator={LaTeX via pandoc}}


\title{Environmental Informatics Using Research Infrastructures and
their Data}
\author{Dr.~Katharyn Duffy, Dr.~Ben Ruddell}
\date{}
\begin{document}
\frontmatter
\maketitle

\renewcommand*\contentsname{Table of contents}
{
\setcounter{tocdepth}{2}
\tableofcontents
}

\mainmatter
\bookmarksetup{startatroot}

\chapter{Environmental Informatics Using Research Infrastructures and
their
Data}\label{environmental-informatics-using-research-infrastructures-and-their-data}

\begin{quote}
2020 Edition
\end{quote}

\bookmarksetup{startatroot}

\chapter*{Preface}\label{preface}
\addcontentsline{toc}{chapter}{Preface}

\markboth{Preface}{Preface}

This textbook provides an introduction to environmental and ecological
informatics in the context of ``big science''- that is, in the context
of research infrastructures and observatories that collect and publish
reams of observations and derived data products about the earth and its
environment. The reader will learn how to make use of environmental
infrastructures' data products.

This textbook introduces a framework of learning outcomes that covers
the broad context of these data products that is necessary to make
proper use of the information:

\begin{itemize}
\tightlist
\item
  The infrastructure's organizational structure and scientific scope;
\item
  Instrumentation, quality control, metadata, data catalogs, API's, and
  data products;
\item
  Key references, key tutorials, and informatics best practices;
\item
  Professional career tracks in informatics;
\end{itemize}

This textbook is intended primarily for graduate students enrolled in
computer science and informatics programs but engaged in studies and
research on environmental and geoscience topics- with a special emphasis
on ecological topics. Advanced undergraduates and other graduate
students enrolled in STEM programs may also be well-served if they have
a strong background in programming and computing. Additionally,
professional scientists may find this textbook useful as a reference and
as a training manual when they encounter the need to make use of the
research infrastructures and data that are directly covered by the
book's content- or find the need to train a junior scientist on the use
of these infrastructures' data.

The scope of the infrastructures and data products is mostly U.S.
focused in this edition, but some of these infrastructures have a global
reach, and the material is almost as useful for students in other
countries as for U.S. students. The textbook can be tackled one unit at
a time as a lab manual within a university course, or- in its intended
application- a standard semester-long three-credit-hour graduate course
should be offered to cover the entire textbook from start to finish.
Digital supplements are provided with examples of successful projects.

Efforts have been made to select activities using data products,
software, and tools that are relatively mature and stable. Even so,
because this textbook covers a rapidly moving field, portions will,
unfortunately, become dated quickly. It is the authors' intent to
release frequent editions that update and expand the material to keep
pace with the rapid development of our field.

Informatics is arguably \emph{the key scientific discipline} of the 21st
century, and research infrastructures are the source of the raw natural
resource fueling the informatics revolution: observational data. Most
21st century scientists and scientific staff will spend their careers
immersed in the data revolution. We sincerely hope that this textbook
provides the launchpad you need for your career or for your next project
in environmental science.

\section*{Acknowledgements}\label{acknowledgements}
\addcontentsline{toc}{section}{Acknowledgements}

\markright{Acknowledgements}

This first version of the textbook was developed to offer INF550, a
graduate course in the School of Informatics Computing and Cyber Systems
at Northern Arizona University in Fall 2020, with funding and leadership
from the National Science Foundation funded National Research
Traineeship ``T3'' option in Ecological and Environmental Informatics
within a PhD program in Informatics and Computing \emph{(NRT-HDR
\#1829075, PI's Ogle, Barber, Richardson, Ruddell, and Sankey)}. Our
research infrastructure partners were critical to the creation of the
material. Our partners at NEON - Battelle deserve special gratitude for
anchoring the project. The opinions expressed are those of the
researchers, and not necessarily the funding agencies. Special thanks to
Megan Jones and Donal O'Leary at NEON-Battelle for their support in
pulling NEON materials.

\subsection*{Key Contributors}\label{key-contributors}
\addcontentsline{toc}{subsection}{Key Contributors}

\emph{Alphabetized by organization, then last name:}

\begin{itemize}
\tightlist
\item
  Alycia Crall, \emph{NEON - Battelle}
\item
  Chris Florian, \emph{NEON - Battelle}
\item
  Megan Jones, \emph{NEON - Battelle}
\item
  Hank Loescher, \emph{NEON - Battelle}
\item
  Paula Mabee, \emph{NEON - Battelle}
\item
  Donal O'Leary, \emph{NEON - Battelle}
\item
  Kate Thibault, \emph{NEON - Battelle}
\item
  Andrew Richardson, \emph{PhenoCam - NAU}
\item
  Bijan Seyednasrollah, \emph{PhenoCam - NAU}
\item
  Theresa Crimmons, \emph{USA-NPN}
\item
  Kathy Gerst, \emph{USA-NPN}
\item
  Lee Marsh, \emph{USA-NPN}
\end{itemize}

This work is licensed under a Creative Commons
Attribution-NonCommercial-ShareAlike 4.0 International License.

\bookmarksetup{startatroot}

\chapter*{Pre-Course Setup: EcoInformatics
Tools}\label{pre-course-setup-ecoinformatics-tools}
\addcontentsline{toc}{chapter}{Pre-Course Setup: EcoInformatics Tools}

\markboth{Pre-Course Setup: EcoInformatics Tools}{Pre-Course Setup:
EcoInformatics Tools}

The purpose of this course is to train you in key ecoinformatics
practices.

Therefore, as an Ecoinformatician you \emph{need} to be able to:

\begin{enumerate}
\def\labelenumi{\arabic{enumi}.}
\item
  Pull data from Application Programming Interfaces (APIs)

  \begin{itemize}
  \tightlist
  \item
    More on this in Chapter 2
  \end{itemize}
\item
  Organize and document your code and data
\item
  Version control your code to avoid disaster and make it reproducible

  \begin{itemize}
  \tightlist
  \item
    For you, your collaborators, and/or the wider community
  \end{itemize}
\item
  Push your code up to public-facing repositories
\item
  Pull others code from public repositories.
\end{enumerate}

More thoughts on the benefits and power of reproducibility
\href{https://github.com/katharynduffy/ECOSS_reproducible_science}{can
be found here}

To be successful, both in this course and in your careers you will need
these skills. This is why \textbf{they are a requirement} for this
course. If you are already using these skills on a daily basis,
fantastic! If you don't feel that you have mastery in the workflows
listed above we have placed lesson links throughout this chapter so that
you can build these skills and be successful in this course.

\section*{Pre-Course Skills \& Setup}\label{pre-course-skills-setup}
\addcontentsline{toc}{section}{Pre-Course Skills \& Setup}

\markright{Pre-Course Skills \& Setup}

For the purpose of this course we will largely be using the following
tools to access, pull, and explore data:

\begin{enumerate}
\def\labelenumi{\arabic{enumi}.}
\tightlist
\item
  R \& Rstudio
\item
  Git, GitHub, \& Atom.io
\item
  Markdown \& Rmarkdown
\end{enumerate}

As such we will need to install and/or update these tools on your
personal computer \emph{before} our first day of class. While we chose R
for this course, nearly all of the packages and data are fully available
and transferable to Python or other languages. If you'd like to brush up
on your R skills I highly recommend Data Carpentry Boot camp's free
\href{http://swcarpentry.github.io/r-novice-gapminder}{R for
Reproducible Scientific Analysis} course.

\subsection*{Installing or Updating R}\label{installing-or-updating-r}
\addcontentsline{toc}{subsection}{Installing or Updating R}

Please check your version of R. You will need R 3.6.0+

How to check your version in R or RStudio if you already have it:

\begin{verbatim}
    > version
                   _                           
    platform       x86_64-apple-darwin15.6.0   
    arch           x86_64                      
    os             darwin15.6.0                
    system         x86_64, darwin15.6.0        
    status                                     
    major          3                           
    minor          5.1                         
    year           2018                        
    month          07                          
    day            02                          
    svn rev        74947                       
    language       R                           
    version.string R version 3.5.1 (2018-07-02)
    nickname       Feather Spray  
\end{verbatim}

If you don't already have R or need to update it
\href{https://cran.rstudio.com/}{do so here.}

\subsection*{Windows R/RStudio Setup}\label{windows-rrstudio-setup}
\addcontentsline{toc}{subsection}{Windows R/RStudio Setup}

After you have downloaded R, run the .exe file that was just downloaded
Go to the
\href{https://www.rstudio.com/products/rstudio/download/\#download}{RStudio
Download page} Under Installers select RStudio X.XX.XXX - e.g.~Windows
Vista/7/8/10 Double click the file to install it Once R and RStudio are
installed, click to open RStudio. If you don't get any error messages
you are set. If there is an error message, you will need to re-install
the program.

\subsection*{Mac R/RStudio Setup}\label{mac-rrstudio-setup}
\addcontentsline{toc}{subsection}{Mac R/RStudio Setup}

After you have downloaded R, double click on the file that was
downloaded and R will install Go to the
\href{https://www.rstudio.com/products/rstudio/download/\#download}{RStudio
Download page} Under Installers select RStudio 1.2.1135 - Mac OS X XX.X
(64-bit) to download it. Once it's downloaded, double click the file to
install it. Once R and RStudio are installed, click to open RStudio. If
you don't get any error messages you are set. If there is an error
message, you will need to re-install the program.

\section*{Linux R/RStudio Setup}\label{linux-rrstudio-setup}
\addcontentsline{toc}{section}{Linux R/RStudio Setup}

\markright{Linux R/RStudio Setup}

R is available through most Linux package managers. You can download the
binary files for your distribution from CRAN. Or you can use your
package manager. e.g.~for Debian/Ubuntu

\begin{Shaded}
\begin{Highlighting}[]
\FunctionTok{sudo}\NormalTok{ apt{-}get install r{-}base}
\end{Highlighting}
\end{Shaded}

and for Fedora

\begin{Shaded}
\begin{Highlighting}[]
\ExtensionTok{run}\NormalTok{ sudo yum install R}
\end{Highlighting}
\end{Shaded}

To install RStudio, go to the
\href{https://www.rstudio.com/products/rstudio/download/\#download}{RStudio
Download page} Under Installers select the version for your
distribution. Once it's downloaded, double click the file to install it
Once R and RStudio are installed, click to open RStudio. If you don't
get any error messages you are set. If there is an error message, you
will need to re-install the program.

\subsection*{Install basic packages for this
course}\label{install-basic-packages-for-this-course}
\addcontentsline{toc}{subsection}{Install basic packages for this
course}

You can run the following script to make sure all the required packages
are properly installed on your computer.

\begin{Shaded}
\begin{Highlighting}[]
\CommentTok{\# list of required packages}
\NormalTok{list.of.packages }\OtherTok{\textless{}{-}} \FunctionTok{c}\NormalTok{(}
  \StringTok{\textquotesingle{}data.table\textquotesingle{}}\NormalTok{,}
  \StringTok{\textquotesingle{}tidyverse\textquotesingle{}}\NormalTok{,}
  \StringTok{\textquotesingle{}jsonlite\textquotesingle{}}\NormalTok{,}
  \StringTok{\textquotesingle{}jpeg\textquotesingle{}}\NormalTok{,}
  \StringTok{\textquotesingle{}png\textquotesingle{}}\NormalTok{,}
  \StringTok{\textquotesingle{}raster\textquotesingle{}}\NormalTok{,}
  \StringTok{\textquotesingle{}rgdal\textquotesingle{}}\NormalTok{,}
  \StringTok{\textquotesingle{}rmarkdown\textquotesingle{}}\NormalTok{, }
  \StringTok{\textquotesingle{}knitr\textquotesingle{}}
\NormalTok{)}

\CommentTok{\# identify new (not installed) packages}
\NormalTok{new.packages }\OtherTok{\textless{}{-}}\NormalTok{ list.of.packages[}\SpecialCharTok{!}\NormalTok{(list.of.packages }\SpecialCharTok{\%in\%} \FunctionTok{installed.packages}\NormalTok{()[,}\StringTok{"Package"}\NormalTok{])]}

\CommentTok{\# install new (not installed) packages}
\ControlFlowTok{if}\NormalTok{(}\FunctionTok{length}\NormalTok{(new.packages)) }
  \FunctionTok{install.packages}\NormalTok{(new.packages, }
                   \AttributeTok{repos=}\StringTok{\textquotesingle{}http://cran.rstudio.com/\textquotesingle{}}\NormalTok{)}

\CommentTok{\# load all of the required libraries}
\FunctionTok{sapply}\NormalTok{(list.of.packages, library, }\AttributeTok{character.only =}\NormalTok{ T)}
\end{Highlighting}
\end{Shaded}

\textbf{Note}: On some operating systems, you may need to install the
Geospatial Data Abstraction Library (GDAL). More information about GDAL
can be found from
\href{https://trac.osgeo.org/gdal/wiki/DownloadingGdalBinaries}{here}.

\section*{Installing and Setting up Git \& Github on Your
Machine}\label{installing-and-setting-up-git-github-on-your-machine}
\addcontentsline{toc}{section}{Installing and Setting up Git \& Github
on Your Machine}

\markright{Installing and Setting up Git \& Github on Your Machine}

For this course you will need: 1. Git installed on your local machine 2.
Very basic bash scripting 3. A linked GitHub account 4. To link RStudio
to git via RStudio or Atom.io

As we will be using these skills constantly, they are a
\emph{pre-requisite} for this course. If you don't yet have these skills
it's okay! You can learn everything that you need to know via the
following freely available resources:

\begin{itemize}
\item
  \href{http://swcarpentry.github.io/shell-novice}{The Unix Shell}
\item
  \href{http://swcarpentry.github.io/git-novice}{Version Control with
  Git}
\item
  \href{https://happygitwithr.com/}{Happy Git with R}
\end{itemize}

If you are learning these skills from scratch I estimate that you will
need to devote \textasciitilde4-6 hours to get set up and comfortable
with the various workflows. Also remember that I have code office hours
every week and that Stack Exchange is your friend.

\section*{Installing Atom}\label{installing-atom}
\addcontentsline{toc}{section}{Installing Atom}

\markright{Installing Atom}

\href{https://atom.io/}{Atom.io} is a powerful and useful text editor
for the following reasons:

\begin{enumerate}
\def\labelenumi{\arabic{enumi}.}
\item
  It is language agnostic
\item
  It fully integrates with git and github + You can use it to
  push/pull/resolve conflicts and write code all in one space.
\end{enumerate}

\section*{Linking RStudio to Git}\label{linking-rstudio-to-git}
\addcontentsline{toc}{section}{Linking RStudio to Git}

\markright{Linking RStudio to Git}

\href{https://happygitwithr.com/rstudio-git-github.html}{Happy Git with
R} has a fantastic tutorial to help you link Rstudio-Git-Github on your
local machine and push/pull from or to public repositories.

\section*{How we will be Conducting this
Course}\label{how-we-will-be-conducting-this-course}
\addcontentsline{toc}{section}{How we will be Conducting this Course}

\markright{How we will be Conducting this Course}

If you find a \textbf{broken link or error} in this course text
\href{https://github.com/katharynduffy/katharynduffy.github.io/issues}{submit
an issue on the course github repository}.

At the end of each chapter you will find a set of \textbf{Exercises}. At
the end of the assigned chapter you will be expected to submit two files
to the course webdrive:

\begin{enumerate}
\def\labelenumi{\arabic{enumi}.}
\item
  An
  \href{https://rstudio.com/wp-content/uploads/2015/02/rmarkdown-cheatsheet.pdf}{RMarkdown
  file} with the naming convention: LASTNAME\_COURSECODE\_Section\#.Rmd,
  and
\item
  A knitted .PDF with the same naming convention:
  LASTNAME\_COURSECODE\_Section\#.pdf
\end{enumerate}

To generate these files you have two options:

\begin{enumerate}
\def\labelenumi{\arabic{enumi}.}
\item
  Click on the pencil and pad logo in the top of this text, copy the
  exercise section code, and drop it into your own .Rmd.
\item
  Git clone our
  \href{https://github.com/katharynduffy/katharynduffy.github.io}{course
  Github Repository}, navigate to the `\_Exercises' folder, and use that
  .Rmd as a template.
\end{enumerate}

\emph{Note: Exercises submitted in any other format, or those missing
questions will not be graded}

To generate your .PDF to upload, in your RMarkdown file simply push the
`Knit' button at the top of your document.

\section*{Exercises:}\label{exercises}
\addcontentsline{toc}{section}{Exercises:}

\markright{Exercises:}

\subsection*{Exercise 0.1: A git
introduction}\label{exercise-0.1-a-git-introduction}
\addcontentsline{toc}{subsection}{Exercise 0.1: A git introduction}

\begin{enumerate}
\def\labelenumi{\arabic{enumi}.}
\tightlist
\item
  Navigate to our
  \href{https://github.com/katharynduffy/Environmental-Informatics-Using-Research-Infrastructures-and-their-Data}{course
  github}
\item
  \texttt{git\ fork} our repo onto your own personal github account.
\item
  \texttt{git\ clone} the repo onto your own personal machine in a place
  that is functional and not temporary \emph{(e.g.~not your downloads
  folder)}.
\end{enumerate}

\begin{Shaded}
\begin{Highlighting}[]
    \CommentTok{\#hints}
    \BuiltInTok{cd} \KeywordTok{\textasciigrave{}}\ExtensionTok{Your/Path/Here}\StringTok{\textquotesingle{}}
\StringTok{    git clone \textquotesingle{}}\ExtensionTok{repo}\NormalTok{ HTTPS}\StringTok{\textquotesingle{}}
\end{Highlighting}
\end{Shaded}

\begin{enumerate}
\def\labelenumi{\arabic{enumi}.}
\setcounter{enumi}{3}
\tightlist
\item
  Add 2-3 sentences introducing yourself in the
  \texttt{\_Course-participants} folder. For example:
\end{enumerate}

\begin{verbatim}
Hi, I'm Dr. Katharyn Duffy.  I have a Ph.D in Earth Science from Northern Arizona University.  Over the past two years I've worked as an open-source software engineer in the PhenoCam lab, and now I'm the coding and lab support for your course.  I really look forward to working with all of you!
\end{verbatim}

\begin{verbatim}
***
Hi, I'm Dr. Katharyn Duffy.  I have a Ph.D in Earth Science from Northern Arizona University.  Over the past two years I've worked as an open-source software engineer in the PhenoCam lab, and now I'm the coding and lab support for your course.  I really look forward to working with all of you!
***
\end{verbatim}

\begin{enumerate}
\def\labelenumi{\arabic{enumi}.}
\setcounter{enumi}{4}
\tightlist
\item
  Submit a \texttt{pull\ request} to add your introduction to our course
  participants folder.
\end{enumerate}

\begin{verbatim}
#hints
git add ...
git commit ...
git status....
git push --set-upstream
git remote -v
git remote add upstream...
\end{verbatim}

\textbf{Note:} \emph{You may complete these either on the command line
or via a program like Atom.io. If you haven't yet made commits to a
remote repository or submitted pull requests please reference the
resources listed above.}

\bookmarksetup{startatroot}

\chapter*{Why `EcoInformatics'?}\label{why-ecoinformatics}
\addcontentsline{toc}{chapter}{Why `EcoInformatics'?}

\markboth{Why `EcoInformatics'?}{Why `EcoInformatics'?}

\emph{Portions of the following introduction were adapted from
\href{https://www.cell.com/trends/ecology-evolution/fulltext/S0169-5347(11)00339-9}{Michener
\& Jones 2012, Trends in Ecology \& Evolution 'Ecoinformatics:
supporting ecology as a data-intensive science}}

Ecology is increasingly becoming a data-intensive science, relying on
massive amounts of data collected by both remote-sensing platforms and
sensor networks embedded in the environment. New observatory networks,
such as the US National Ecological Observatory Network (NEON), provide
research platforms that enable scientists to examine phenomena across
diverse ecosystem types through access to thousands of sensors
collecting diverse environmental observations. These networks spatially
and temporaly overlap with a number of other networks and
infrastructures ranging from remote sensing, to citizen science, and so
on.

It has been argued that \textbf{data-intensive science represents the
fourth scientific paradigm} following the empirical
\emph{(i.e.~description of natural phenomena)}, theoretical
\emph{(e.g.~modeling and generalization)} and computational
\emph{(e.g.~simulation)} scientific approaches, and comprises an
approach for unifying theory, experimentation and simulation.

Ecologists increasingly address questions at broader scales that have
both scientific and societal relevance. For example, the
\href{https://academic.oup.com/bioscience/article/61/4/290/324883}{40
top priorities for science that can inform conservation and management
policy in the USA} rely principally on a sound foundation of ecological
research, and the ability to scale knowledge and inter-connect data.

\includegraphics[width=6.34in,height=\textheight,keepaspectratio]{./images/interconnect_data.png}

\begin{quote}
Continental-scale patterns and dynamics result from climate and people
as broad-scale drivers interacting with finer-scale vectors that
redistribute materials within and among linked terrestrial and aquatic
systems. Climate and land-use change interact with patterns and
processes at multiple, finer scales (blue arrows). (a) These drivers can
influence broad-scale patterns directly, and these constraints may act
to overwhelm heterogeneity and processes at (b) mesoscales and at (c)
the finer scale of local sites. Broad-scale drivers can also exert an
indirect impact on broad-scale patterns through their interactions with
disturbances, including (d) the spread of invasive species, (e)
pattern--process relationships at meso-scales, or (f) at finer scales
within a site. Connectivity imparted by the transfer of materials occurs
both at (g) the meso-scale and at (h) finer scales within sites where
terrestrial and aquatic systems are connected. These dynamics at fine
scales can propagate to influence larger spatial extents (red arrows).
Feedbacks occur throughout the system. The term ``drivers'' refers to
both forcing functions that are part of the system and to external
drivers.
\href{https://esajournals.onlinelibrary.wiley.com/doi/pdf/10.1890/070098}{Peters
et al., 2008}
\end{quote}

\textbf{Ecology is also affected by changes that are occurring
throughout science as a whole.}

In particular, scientists, professional societies and research sponsors
are recognizing the value of data as a product of the scientific
enterprise and placing increased emphasis on data stewardship, data
sharing, openness and supporting study repeatability.

\includegraphics[width=5.83in,height=\textheight,keepaspectratio]{./images/intro_data.png}

\begin{quote}
Data on ecological and environmental systems are (A) acquired, checked
for quality, documented using an acquisition workflow, and then both the
raw and derived data products are versioned and deposited in the DataONE
federated data archive (red dashed arrows). Researchers discover and
access data from the federation and then (B) integrate and process the
data in an analysis workflow, resulting in derived data products,
visualizations, and scholarly papers that are in turn archived in the
data federation (red dashed arrows). Other researchers directly cite any
of the versioned data, workflows, and visualizations that are archived
in the DataONE federation.
\href{https://science.sciencemag.org/content/sci/331/6018/703.full.pdf?casa_token=dDumP2G8NpEAAAAA:8Zrx01SOGOQK79my5SKJdsy5buDQAOd_wVclvlq2eFJZ0MWpeUlHQ7zx3ygkfoQOXSyRLhgx0FDO9V4}{Richman
et al., 2011}
\end{quote}

\textbf{The changes that are occurring in ecology create challenges with
respect to acquiring, managing and analyzing the large volumes of data
that are collected by scientists worldwide.}

One challenge that is particularly daunting lies in dealing with the
scope of ecology and the enormous variability in scales that is
encountered, spanning microbial community dynamics, communities of
organisms inhabiting a single plant or square meter, and ecological
processes occurring at the scale of the continent and biosphere. The
diversity in scales studied and the ways in which studies are carried
out results in large numbers of
\href{https://muse.jhu.edu/article/262029/pdf?casa_token=8ysDn-0aPmAAAAAA:kqtGap8fEbZOvCZE_tcp4sjof-RwcLsiNvk0kk4fpx90jv4npnStRQAPQOb4jb9UC-JINBLlCEw}{small,
idiosyncratic data sets} that accumulate from the thousands of
scientists that collect relevant biological, ecological and
environmental data.

\includegraphics[width=8.17in,height=\textheight,keepaspectratio]{./images/data-ontology.png}

\begin{quote}
A proposed high-level architecture for ecological and environmental data
management is shown consisting of three primary levels. Data stored
within distributed data repositories (a) is mediated by standard
metadata and ontologies (b) to power software tools used by scientists
and data managers (c). Software applications use community-endorsed
ontologies and metadata standards from the middle level to provide tools
that are more effective for publishing, querying, integrating and
analyzing data. Ontologies are separated into framework ontologies and
domain-specific extensions, enabling contributions from multiple
research groups, disciplines and individuals. Cross-disciplinary data
are maintained in local repositories, but made accessible to the broader
research community through distributed systems based on shared, open
protocols (such as Metacat). Example repositories include the LTER
network, National Ecological Observatory Network, United States
Geographical Survey and SEEK's EarthGrid.
\href{https://www.sciencedirect.com/science/article/pii/S0169534708000384?casa_token=b8GX-9jkuLwAAAAA:nQtJlhAUGyZdcokpQLEtho3OfVeDUKhENueO0k8H7dWPmPmXOUFYhGR6sKNtT04rEgJNR_HNlqg}{Madin
et al. 2008,}
\end{quote}

\textbf{Ecoinformatics is a framework} that enables scientists to
generate new knowledge through innovative tools and approaches for:

\begin{itemize}
\tightlist
\item
  discovering,
\item
  managing,
\item
  integrating,
\item
  analyzing,
\item
  visualizing,
\item
  and preserving
\end{itemize}

relevant biological, environmental, and socioeconomic data and
information. Many ecoinformatics solutions have been developed
\emph{over the past decade}, increasing scientists' efficiency and
supporting faster and easier data discovery, integration and analysis;
however, many challenges remain, especially in relation to installing
ecoinformatics practices into mainstream research and education.

And that, course participants, is why we are here.

\section*{The Framework of this
Course}\label{the-framework-of-this-course}
\addcontentsline{toc}{section}{The Framework of this Course}

\markright{The Framework of this Course}

Over the duration of this course we will survey a wide array of
observation platforms and networks and build hands-on experience with
the framework of Ecoinformatics. For coherance we will cover the
following overarching themes:

\begin{quote}
\begin{enumerate}
\def\labelenumi{\arabic{enumi}.}
\tightlist
\item
  Each network's \textbf{mission and design}
\item
  Each network's \textbf{spatial design}
\end{enumerate}

\begin{itemize}
\tightlist
\item
  e.g.~opprtunistic vs.~planned, citizen science vs.~orbital sensors
\end{itemize}

\begin{enumerate}
\def\labelenumi{\arabic{enumi}.}
\setcounter{enumi}{2}
\tightlist
\item
  The \textbf{types of data} that stream from each network
\end{enumerate}

\begin{itemize}
\tightlist
\item
  e.g.~sensors, derived products, metadata
\end{itemize}

\begin{enumerate}
\def\labelenumi{\arabic{enumi}.}
\setcounter{enumi}{3}
\tightlist
\item
  \textbf{How to access that data}
\end{enumerate}

\begin{itemize}
\tightlist
\item
  e.g.~APIs, landing pages, r packages etc.
\end{itemize}

\begin{enumerate}
\def\labelenumi{\arabic{enumi}.}
\setcounter{enumi}{4}
\tightlist
\item
  Opportunities to interact with or contribute to each network
\end{enumerate}

\begin{itemize}
\tightlist
\item
  e.g.~RFP's coming down the pipeline, internships, and post-doctoral
  scholar programs.
\end{itemize}
\end{quote}

At the conclusion of each network's section you will be asked to write a
1-page summary reviewing the above framework for each network, and
highlight how it potentially aligns with your own research. These series
of 1-page summaries will then culminate into a final presentation where
you propose to derive your own data product for your own research
\emph{touching upon multiple networks} and \emph{accounting for
differences in spatial footprints, frequency of observations, and
important data cross-walks}.

\section*{Final Course Project: Proposed Derived Data
Product}\label{final-course-project-proposed-derived-data-product}
\addcontentsline{toc}{section}{Final Course Project: Proposed Derived
Data Product}

\markright{Final Course Project: Proposed Derived Data Product}

For your final project, you will present a \textbf{4-6 minute
IGNITE-style derived data product pitch}, followed by \emph{2-3 minutes
of questions from your audience} (which will include members from the
infrastrures we've covered). Think of this project as your \emph{`sales
pitch'} to the research infrastructure whose data you are using, and/or
the scientific community as a whole. In the IGNITE theme of
\emph{`Enlighten us but make it quick'}, you will construct a series of
slides that auto-advance every 30 seconds. Specific instructions for the
content of each slide are below.

Ideally, this final presentation will feed upon a number of the
`culmination write-ups' you have conducted over the course of the
semester. Ideally, this derived data product will utilize data from a
number of sources, either covered within this course or external to it.
Ideally, it will also convince your audience that your idea is
\textbf{novel, useful,} \emph{and} \textbf{possible}.

In order to complete this presentation, you will need to have worked
with the various data products you propose, have an in-depth
understanding of them, and their challenges, along with original, clean,
high-level summary graphics. Further, giving an IGNITE-style
presentation takes practice. IGNITE-style presentations are powerful, as
they keep \emph{you} moving forward, and give your audience a
\emph{high-level understanding} of your topic. We fully recommend
rehearsing your presentation \textbf{many times} before giving it live
and recording yourself to learn how you can improve.

Here's an example (of an even faster) ignite talk from one of your
book's authors:

\textbf{In your derived data product pitch you will cover these themes:}

\begin{enumerate}
\def\labelenumi{\arabic{enumi}.}
\item
  The need for the derived data that you are proposing to produce.
\item
  What data you will use to derive this product, including the
  justification for this exact data.
\item
  The processing pipeline for this product, along with estimates for a
  timeline.
\item
  Potential hurdles you will have to overcome.
\item
  How this product will serve the infrastructure and/or the scientific
  community.
\end{enumerate}

\textbf{Specific slide criteria are as follows:}

\textbf{Slide 1:} Title, authors (including contacts at infrastructures
covered if applicable)

\textbf{Slide 2:} Justification for the derived data product; the gap or
need that it fills

\textbf{Slides 3-x:}

1 slide per data product used including:

\begin{itemize}
\item
  The exact data product (e.g.~NEON data product id and full title)
\item
  A 1 sentence summary of the data product and its justification for
  this purpose
\item
  An original, clean, polished high-level plot, gif or .mp4 of the data
\end{itemize}

\textbf{Slide x + 1:} A high-level workflow diagram of the processing
pipeline \emph{E.g.: Original data and how you pull it in (API,
r-package etc)}

\begin{itemize}
\item
  Filtering process using QA/QC or metadata
\item
  Orthorectification in time or space
\end{itemize}

Example generated using \href{https://app.diagrams.net/}{draw.io}:

\includegraphics[width=3.28in,height=\textheight,keepaspectratio]{./images/APIS_usecase2_v2 (4).png}

\textbf{Slide x + 2:} A clean plot of all of the data you mentioned
together, and/or the derived data product itself with a 1 sentence
summary

Example:

\includegraphics[width=35.01in,height=\textheight,keepaspectratio]{./images/mased_APIS_2.png}

\textbf{Slide x+3:} Summary: Circle back on how this derived data
product serves your research, the infrastructure, and the wider science
community (no more than 10 words, suggestion: graphics or bullet points)

\textbf{Slide x + 4:} Data citations for all data used in proposed
derived data product

\textbf{An example slide deck with specific ideas}
\href{https://docs.google.com/presentation/d/16vRomNXmdyGrKQ_bdsGaOfbbU-CW_ZGi5-MPhzDBclY/edit?usp=sharing}{\textbf{can
be found here}}

\textbf{The rubric for your final presentation grade is as follows:}

Presentation meets all requirements and criteria: 60\%

Aesthetics and craft of presentation: 10\%

Live presentation of materials: 30\%


\backmatter


\end{document}
